% Generated draft: requires mathematical and source review.
\documentclass[11pt]{article}
\usepackage[margin=25mm]{geometry}
\usepackage{fontspec}
\setmainfont{Latin Modern Roman}
\usepackage{amsmath,amssymb,mathtools,unicode-math}
\setmathfont{Latin Modern Math}
\usepackage{enumitem,needspace,xurl,hyperref}
\setlength{\parindent}{0pt}
\setlength{\parskip}{5pt}
\setlength{\emergencystretch}{3em}
\newcommand{\N}{\mathbb{N}}
\newcommand{\Z}{\mathbb{Z}}
\newcommand{\Q}{\mathbb{Q}}
\newcommand{\R}{\mathbb{R}}
\newcommand{\C}{\mathbb{C}}
\newcommand{\F}{\mathbb{F}}
\newcommand{\A}{\mathbb{A}}
\newcommand{\PP}{\mathbb{P}}
\newcommand{\E}{\mathbb{E}}
\newcommand{\eps}{\varepsilon}
\newcommand{\dd}{\,\mathrm{d}}
\begin{document}
\section*{UnsolvedMath ID: 801155}
\textbf{Source problem code:} catalogue entry
\subsection{Definitions and mathematical statement}
Let \(M\) be a smooth real manifold of dimension \(n\), and let \(L\colon TM\to TM\) be a smooth field of endomorphisms. Its Nijenhuis torsion is
\[
N_L(X,Y)=[LX,LY]-L[LX,Y]-L[X,LY]+L^2[X,Y]
\]
for smooth vector fields \(X,Y\). The operator \(L\) is Nijenhuis if \(N_L=0\). It is \(\mathrm{gl}\)-regular if, at every point \(p\), the minimal polynomial of \(L_p\) has degree \(n\). In dimension two this means that \(L_p\) is not a scalar multiple of the identity at any point.

A standard interpretation of a metric geodesically compatible with \(L\) is a smooth nondegenerate symmetric tensor \(g\) such that \(L\) is self-adjoint with respect to \(g\) and
\[
(\nabla_XL)Y
=\frac12\left(\dd(\operatorname{tr}L)(Y)X
+g(X,Y)\operatorname{grad}_g(\operatorname{tr}L)\right),
\]
where \(\nabla\) is the Levi--Civita connection of \(g\), and \(g(\operatorname{grad}_g h,X)=\dd h(X)\). This convention is provided to make a possible interpretation explicit; the supplied excerpt does not define geodesic compatibility or specify the allowed metric signature, and its intended convention requires verification.

\textbf{Requested construction, Problem 8.} For each \(\mathrm{gl}\)-regular Nijenhuis operator \(L\) in dimension two, construct explicitly a geodesically compatible metric. The source does not specify whether the construction is local near every point or global on the given manifold. Consequently, a global existence assertion, a choice of signature, and any additional invertibility requirement on \(L\) are not imposed here. The real smooth setting above is a conventional interpretation rather than an explicit hypothesis supplied by the excerpt.

\textbf{Additional requested generalization.} In arbitrary coordinates \(x^1,\ldots,x^n\), assume that a conservation law is supplied as an explicit function \(f(x^1,\ldots,x^n)\). Generalize the formula referred to in the source as the above-mentioned formula from \texttt{[nijappl5]} to these coordinates. The supplied record contains neither that formula nor the bibliographic identification of \texttt{[nijappl5]}, and it does not define the conservation-law condition on \(f\). This part therefore specifies a requested generalization but does not determine a precise formula or a complete set of hypotheses. Its relationship to the dimension-two construction also remains unspecified.
\subsection{Short English statement}
For every two-dimensional \(\mathrm{gl}\)-regular Nijenhuis operator, explicitly construct a geodesically compatible metric. The record also asks for a coordinate generalization of an omitted formula using a given conservation-law function. Metric signature, local versus global scope, and the missing formula and conservation-law definition require clarification from the source. The supplied registry labels the problem open; this status is not independently verified.
\subsection{Sources}
\begin{itemize}
\item Local UnsolvedMath registry, permanent ID \(801155\), exact source problem code \texttt{catalogue entry}, title \emph{Nijenhuis Geometry and Integrable Systems --- Problem 8}. The supplied record attributes the problem to Bolsinov, Konyaev, and Matveev, with proposed year 2024, and labels its status \texttt{open}.
\item Bolsinov, Konyaev, Matveev, \emph{Research Problems on Nijenhuis Geometry and Integrable Systems} (2024), Problem 8. Supplied source URL: \url{https://arxiv.org/abs/2410.04276}. The provenance records extraction as \texttt{arxiv\_latex} and access on 2026-10-16. These are supplied provenance statements; no independent source visit or literature verification is claimed.
\item The statement refers to \texttt{[nijappl5]}, but the supplied record includes neither its bibliographic entry nor the referenced formula.
\end{itemize}
\end{document}
